% problem_id: S001-etale-cohomology-lemma-points-small-etale-site
% source_id: S001 (Stacks Project)
% source_file: etale-cohomology.tex
% statement_locator: etale-cohomology.tex lines 3609--3626
% proof_locator: etale-cohomology.tex lines 3628--3790
% env: lemma
% id_note: label 跨文件重名 ⇒ id 用文件限定形式
% pairing: tight (verified)
% status: complete (statement + author proof verbatim + reason + locators)
%
\begin{lemma}
\label{lemma-points-small-etale-site}
Let $S$ be a scheme.
\begin{enumerate}
\item Let $p$ be a point of the small \'etale site
$S_\etale$ of $S$ given by a functor
$u : S_\etale \to \textit{Sets}$.
Then there exists a geometric point $\overline{s}$ of $S$ such that
$p$ is isomorphic to the point of $S_\etale$ associated to
$\overline{s}$ in
Lemma \ref{lemma-stalk-gives-point}.
\item Let $p : \Sh(pt) \to \Sh(S_\etale)$ be a point
of the small \'etale topos of $S$. Then $p$ comes from a geometric point
of $S$, i.e., the stalk functor $\mathcal{F} \mapsto \mathcal{F}_p$
is isomorphic to a stalk functor as defined in
Definition \ref{definition-stalk}.
\end{enumerate}
\end{lemma}

\begin{proof}
By
Sites, Lemma \ref{sites-lemma-point-site-topos}
there is a one to one correspondence between points of the site and points
of the associated topos, hence it suffices to prove (1).
By
Sites, Proposition \ref{sites-proposition-point-limits}
the functor $u$ has the following properties:
(a) $u(S) = \{*\}$, (b) $u(U \times_V W) = u(U) \times_{u(V)} u(W)$, and
(c) if $\{U_i \to U\}$ is an \'etale covering, then
$\coprod u(U_i) \to u(U)$ is surjective.
In particular, if $U' \subset U$ is an open subscheme,
then $u(U') \subset u(U)$. Moreover, by
Sites, Lemma \ref{sites-lemma-point-site-topos}
we can write $u(U) = p^{-1}(h_U^\#)$, in other words $u(U)$ is the
stalk of the representable sheaf $h_U$. If
$U = V \amalg W$, then we see that $h_U = (h_V \amalg h_W)^\#$ and we get
$u(U) = u(V) \amalg u(W)$ since $p^{-1}$ is exact.

\medskip\noindent
Consider the restriction of $u$ to $S_{Zar}$. By
Sites, Examples \ref{sites-example-point-topological} and
\ref{sites-example-point-topology}
there exists a unique point $s \in S$ such that for $S' \subset S$ open we
have $u(S') = \{*\}$ if $s \in S'$ and $u(S') = \emptyset$ if $s \not \in S'$.
Note that if $\varphi : U \to S$ is an object of $S_\etale$ then
$\varphi(U) \subset S$ is open (see
Proposition \ref{proposition-etale-morphisms})
and $\{U \to \varphi(U)\}$ is an \'etale covering. Hence we conclude that
$u(U) = \emptyset \Leftrightarrow s \not \in \varphi(U)$.

\medskip\noindent
Pick a geometric point $\overline{s} : \overline{s} \to S$ lying over $s$, see
Definition \ref{definition-geometric-point}
for customary abuse of notation. Suppose that $\varphi : U \to S$ is an object
of $S_\etale$ with $U$ affine. Note that $\varphi$ is separated, and
that the fibre $U_s$ of $\varphi$ over $s$ is an affine scheme over
$\Spec(\kappa(s))$ which is the spectrum of a finite product of
finite separable extensions $k_i$ of $\kappa(s)$. Hence we may apply
\'Etale Morphisms, Lemma \ref{etale-lemma-etale-etale-local-technical}
to get an \'etale neighbourhood $(V, \overline{v})$ of $(S, \overline{s})$
such that
$$
U \times_S V = U_1 \amalg \ldots \amalg U_n \amalg W
$$
with $U_i \to V$ an isomorphism and $W$ having no point lying over
$\overline{v}$. Thus we conclude that
$$
u(U) \times u(V) =
u(U \times_S V) =
u(U_1) \amalg \ldots \amalg u(U_n) \amalg u(W)
$$
and of course also $u(U_i) = u(V)$. After shrinking $V$ a bit we can
assume that $V$ has exactly one point lying over $s$, and hence $W$ has no
point lying over $s$. By the above this then gives $u(W) = \emptyset$.
Hence we obtain
$$
u(U) \times u(V) =
u(U_1) \amalg \ldots \amalg u(U_n) =
\coprod\nolimits_{i = 1, \ldots, n} u(V)
$$
where the equality signs are compatible with the maps to the set $u(V)$.
Note that $u(V) \not = \emptyset$ as $s$ is in the image of $V \to S$.
In particular, we see that in this situation $u(U)$ is a finite
set with $n$ element, namely the left hand side has fibre $u(U)$ over
any element of $u(V)$ and the right hand side has fibre of cardinality $n$.

\medskip\noindent
Consider the limit
$$
\lim_{(V, \overline{v})} u(V)
$$
over the category of \'etale neighbourhoods $(V, \overline{v})$ of
$\overline{s}$. It is clear that we get the same value when taking
the limit over the subcategory of $(V, \overline{v})$ with $V$ affine.
By the previous paragraph (applied with the roles of $V$ and $U$ switched)
we see that in this case $u(V)$ is always a finite nonempty set.
Moreover, the limit is cofiltered, see
Lemma \ref{lemma-cofinal-etale}.
Hence by
Categories, Section \ref{categories-section-codirected-limits}
the limit is nonempty. Pick an element $x$ from this limit.
This means we obtain a $x_{V, \overline{v}} \in u(V)$ for
every \'etale neighbourhood $(V, \overline{v})$ of $(S, \overline{s})$
such that for every morphism of \'etale neighbourhoods
$\varphi : (V', \overline{v}') \to (V, \overline{v})$ we have
$u(\varphi)(x_{V', \overline{v}'}) = x_{V, \overline{v}}$.

\medskip\noindent
We will use the choice of $x$ to construct a functorial bijective map
$$
c : |U_{\overline{s}}| \longrightarrow u(U)
$$
for $U \in \Ob(S_\etale)$ which will conclude the proof. See
Lemma \ref{lemma-stalk-gives-point}
and its proof for a description of $|U_{\overline{s}}|$.
First we claim that it suffices to construct the map for $U$ affine.
We omit the proof of this claim.
Assume $U \to S$ in $S_\etale$ with $U$ affine, and let
$\overline{u} : \overline{s} \to U$ be an element of $|U_{\overline{s}}|$.
Choose a $(V, \overline{v})$ such that $U \times_S V$ decomposes
as in the third paragraph of the proof.
Then the pair $(\overline{u}, \overline{v})$ gives a geometric point of
$U \times_S V$ lying over $\overline{v}$ and determines one of the
components $U_i$ of $U \times_S V$. More precisely, there exists
a section $\sigma : V \to U \times_S V$ of the projection $\text{pr}_U$
such that $(\overline{u}, \overline{v}) = \sigma \circ \overline{v}$. Set
$c(\overline{u}) = u(\text{pr}_U)(u(\sigma)(x_{V, \overline{v}})) \in u(U)$.
We have to check this is independent of the choice of $(V, \overline{v})$. By
Lemma \ref{lemma-cofinal-etale}
the category of \'etale neighbourhoods is cofiltered.
Hence it suffice to show
that given a morphism of \'etale neighbourhood
$\varphi : (V', \overline{v}') \to (V, \overline{v})$ and a choice of a
section $\sigma' : V' \to U \times_S V'$ of the projection such that
$(\overline{u}, \overline{v'}) = \sigma' \circ \overline{v}'$
we have $u(\sigma')(x_{V', \overline{v}'}) = u(\sigma)(x_{V, \overline{v}})$.
Consider the diagram
$$
\xymatrix{
V' \ar[d]^{\sigma'} \ar[r]_\varphi & V \ar[d]^\sigma \\
U \times_S V' \ar[r]^{1 \times \varphi} &
U \times_S V
}
$$
Now, it may not be the case that this diagram commutes. The reason is
that the schemes $V'$ and $V$ may not be connected, and hence
the decompositions used to construct $\sigma'$ and $\sigma$ above may
not be unique. But we do know that
$\sigma \circ \varphi \circ \overline{v}' =
(1 \times \varphi) \circ \sigma' \circ \overline{v}'$
by construction. Hence, since $U \times_S V$ is \'etale over $S$,
there exists an open neighbourhood
$V'' \subset V'$ of $\overline{v'}$ such that the diagram does
commute when restricted to $V''$, see
Morphisms, Lemma \ref{morphisms-lemma-value-at-one-point}.
This means we may extend the diagram above to
$$
\xymatrix{
V'' \ar[r] \ar[d]^{\sigma'|_{V''}} &
V' \ar[d]^{\sigma'} \ar[r]_\varphi &
V \ar[d]^\sigma \\
U \times_S V'' \ar[r] &
U \times_S V' \ar[r]^{1 \times \varphi} &
U \times_S V
}
$$
such that the left square and the outer rectangle commute.
Since $u$ is a functor this implies that
$x_{V'', \overline{v}'}$ maps to the same element in
$u(U \times_S V)$ no matter which route we take through the
diagram. On the other hand, it maps to the elements
$x_{V', \overline{v}'}$ and $x_{V, \overline{v}}$ in
$u(V')$ and $u(V)$. This implies the desired equality
$u(\sigma')(x_{V', \overline{v}'}) = u(\sigma)(x_{V, \overline{v}})$.

\medskip\noindent
In a similar manner one proves that the construction
$c : |U_{\overline{s}}| \to u(U)$ is functorial in $U$;
details omitted. And finally, by the results of the
third paragraph it is clear that the map $c$ is bijective
which ends the proof of the lemma.
\end{proof}
